% !TEX TS-program = pdflatex
\documentclass[10pt,a4paper]{altacv}

% Marges équilibrées pour 1 page complète
\geometry{left=1.2cm,right=1.2cm,top=1.0cm,bottom=0.8cm,footskip=0pt}

\usepackage[T1]{fontenc}
\usepackage[utf8]{inputenc}
\usepackage[default]{lato}
\usepackage{microtype}
\DisableLigatures{encoding=*, family=*} % ATS-fix: Disables fi/fl/ffi ligatures in pdfLaTeX

\definecolor{Navy}{HTML}{0A2540}
\definecolor{SlateGrey}{HTML}{2E2E2E}
\definecolor{LightGrey}{HTML}{555555}
\colorlet{heading}{Navy}
\colorlet{accent}{Navy}
\colorlet{emphasis}{SlateGrey}
\colorlet{body}{LightGrey}

\renewcommand{\itemmarker}{{\small\textbullet}}
\renewcommand{\ratingmarker}{\faCircle}

\usepackage[hidelinks]{hyperref}
\hypersetup{
  pdftitle={Curriculum Vitae -- Achraf Izelmadene},
  pdfauthor={Achraf Izelmadene},
  pdfsubject={Ingénieur Génie Électrique et Systèmes Intelligents},
  pdfkeywords={génie électrique, IoT, automatisation, Edge AI, Embedded Systems, stage}
}

\photo{2.2cm}{assets/imageee.png}

% Style compact pour les sections
\renewcommand{\cvsection}[2][]{%
  \vspace{4pt}%
  {\color{heading}\Large\bfseries #2 \hspace{0.2em}\leaders\vrule height 3pt depth -1.5pt\hfill\mbox{}}\par\vspace{3pt}
}

\renewcommand{\cvevent}[4]{%
  {\normalsize\bfseries\color{accent}#1\par}\vspace{1pt}
  \ifstrequal{#2}{}{}{%
    {\small\itshape\color{emphasis}#2\ifstrequal{#4}{}{}{ - #4}}%
  }%
  ~\hfill
  \ifstrequal{#3}{}{}{{\small\color{emphasis}#3}}%
  \par\vspace{2pt}
}

% Header fusionné avec À Propos
\makeatletter
\newcommand{\target}[1]{\def\@target{#1}}
\newcommand{\apropos}[1]{\def\@apropos{#1}}

\renewcommand{\makecvheader}{%
  \begingroup
    \ifdef{\@photodiameter}{%
    \begin{minipage}[c]{\@photodiameter}
    \if@normalphoto
      \includegraphics[width=\linewidth]{\@photo}
    \else
      \begin{tikzpicture}
        \clip (0,0) circle (0.5\linewidth);
        \node at (0,-0.08\linewidth) {\includegraphics[width=1.05\linewidth]{\@photo}};
      \end{tikzpicture}
    \fi%
    \end{minipage}\hspace{1.2em}%
    \begin{minipage}[c]{\dimexpr\linewidth-\@photodiameter-1.5em}}{}%
    \raggedright\color{emphasis}%
    {\LARGE\bfseries\MakeUppercase{\@name}\par}\vspace{1.5pt}
    {\normalsize\bfseries\color{accent}\@tagline\par}\vspace{3pt}
    \ifdef{\@target}{{\small\textbf{\@target}\par}\vspace{2pt}}{}%
    \ifdef{\@apropos}{{\small\color{emphasis}\@apropos\par}\vspace{3pt}}{}%
    \vspace{2.5pt}%
    {\footnotesize\color{emphasis}\@personalinfo\par}
    \ifdef{\@photodiameter}{%
    \end{minipage}\par}{}%
  \endgroup\vspace{1pt}
}
\makeatother

\begin{document}

\name{Izelmadene Achraf}
\tagline{Élève Ingénieur en Génie Électrique et Systèmes Intelligents (GESI)}

% =====================================================================
% À PROPOS (profil / accroche du CV)
% AGENT INSTRUCTION: \target = objectif de stage (1 ligne), \apropos = pitch (2-3 lignes).
% Ces deux champs s'affichent dans l'en-tête, à côté de la photo.
% Version généraliste GESI : adapter le pitch à l'offre si besoin (sources: kb/).
% =====================================================================
\target{À la recherche d’un stage de fin d’études (PFE) de 6 mois dès février 2027 en génie électrique, automatisme ou systèmes intelligents.}
\apropos{Double compétence en génie électrique (installations HTA/BT, électronique de puissance, automatismes) et en systèmes intelligents (systèmes embarqués, IoT, IA). Fort d’expériences en bureau d’études et en unité de production, je souhaite contribuer à des projets alliant énergie, industrie et innovation.}

\personalinfo{%
  \href{mailto:achraf.izelmadene@gmail.com}{achraf.izelmadene@gmail.com}%
  \hspace{0.9em}\textcolor{body!60}{|}\hspace{0.9em}%
  +212699098736%
  \hspace{0.9em}\textcolor{body!60}{|}\hspace{0.9em}%
  Marrakech, Maroc%
  \hspace{0.9em}\textcolor{body!60}{|}\hspace{0.9em}%
  Permis B%
}

\makecvheader

\cvsection{FORMATION}
\cvevent{Étudiant en cycle d’ingénieur en Génie Électrique et Systèmes Intelligents}{École Marocaine des Sciences de l’Ingénieur - Marrakech}{2025 - Présent}{}
{\footnotesize\color{emphasis}Conception des installations électriques, Électronique de puissance, Réseaux électriques \& Sécurité, Énergies renouvelables \& Smart Grids, Automatique \& Régulation industrielle, Systèmes embarqués, Smart capteurs, Machine Learning \& Deep Learning, Supervision \& Digital Twin, Robotique industrielle, Gestion de projet...}\par\vspace{2.5pt}

\cvevent{Licence SMP, option Automatique}{Faculté des Sciences Semlalia - Marrakech}{2020 - 2024}{}
{\footnotesize\color{emphasis}Électronique analogique, Électromagnétisme, Asservissements linéaires, Automatismes et capteurs, Traitement du signal, Régulation \& logique floue...}\par

\cvsection{EXPÉRIENCE}
\cvevent{Stagiaire Bureau d’Études}{Core Engineering Solutions}{juil. - août 2026}{Berrchid, Casablanca}
\begin{itemize}[leftmargin=*, label=\textbullet, nosep, itemsep=1.5pt]
  \item Dimensionnement d’un poste HTA/BT et intégration d’une installation solaire PV en autoconsommation.
  \item Réalisation de schémas unifilaires, bilans de puissance et notes de calcul (AutoCAD, Caneco BT, PVsyst).
\end{itemize}
\vspace{2pt}
\cvevent{Stagiaire en Unité de Production}{Menara Prefa}{juil. 2024}{Marrakech}
\begin{itemize}[leftmargin=*, label=\textbullet, nosep, itemsep=1.5pt]
  \item Migration d’un automate industriel vers Siemens S7 sous TIA Portal, suivi de maintenance préventive/corrective.
  \item Diagnostic des pannes électriques et participation à l’amélioration continue des processus industriels.
\end{itemize}

\cvsection{PROJETS}
% AGENT INSTRUCTION: Projets choisis selon les étoiles de kb/projects.md (5★, 4★, 4★).
{\normalsize\bfseries\color{accent} Precision Agriculture IoT Ecosystem (Arduino, LoRa, ESP32 \& TFLite)}\par\vspace{1.5pt}
\begin{itemize}[leftmargin=*, label=\textbullet, nosep, itemsep=1.5pt]
  \item Irrigation intelligente : nœuds capteurs (Arduino Nano + LoRa) et unité centrale de décision (ESP32 + logique floue).
  \item Dashboard de monitoring temps réel et modèle TFLite sur smartphone pour la détection de maladies foliaires.
\end{itemize}
\vspace{2pt}
{\normalsize\bfseries\color{accent} Maintenance Prédictive 4.0 par AutoML (TPOT vs. Auto-WEKA)}\par\vspace{1.5pt}
\begin{itemize}[leftmargin=*, label=\textbullet, nosep, itemsep=1.5pt]
  \item Benchmark génétique vs. bayésien sur données IoT (vibrations, température, pression) fortement déséquilibrées.
  \item Pipelines ML automatisés (F1-score, Recall) complétés par une étude CND par analyse d’huile industrielle.
\end{itemize}
\vspace{2pt}
{\normalsize\bfseries\color{accent} Détection d’anomalies Smartgrid (Clustering)}\par\vspace{1.5pt}
\begin{itemize}[leftmargin=*, label=\textbullet, nosep, itemsep=1.5pt]
  \item Comparaison de KMeans, DBSCAN et GMM sur des données IoT de réseaux électriques (Silhouette, Precision, Recall).
  \item Développement d’un dashboard interactif de visualisation des clusters pour le monitoring d’anomalies.
\end{itemize}

\cvsection{COMPÉTENCES}
\begin{tabular}{@{}ll@{}}
\textbf{Ingénierie \& Outils :} & AutoCAD, Caneco BT, PVsyst, Dialux EVO, KiCad, TIA Portal, ISIS Proteus, SolidWorks \\[1.5pt]
\textbf{Systèmes Embarqués \& IoT :} & ESP32, Arduino, LoRa, Modbus RTU, Capteurs, MATLAB / Simulink \\[1.5pt]
\textbf{Programmation \& IA :} & Python, C/C++, Scikit-learn, OpenCV, TensorFlow Lite, AutoML, MLOps \\[1.5pt]
\textbf{Compétences transversales :} & Gestion de projet, Esprit d'équipe, Communication, Entrepreneuriat \& innovation
\end{tabular}

\cvsection{CERTIFICATIONS}
\begin{itemize}[leftmargin=*, label=\textbullet, nosep, itemsep=1.5pt]
  \item \textbf{Foundations of IoT Systems and Industrial Automation}~\hfill {\color{emphasis}LearnQuest (Coursera)}
  \item \textbf{Sensors and Sensor Circuit Design}~\hfill {\color{emphasis}University of Colorado Boulder (Coursera)}
  \item \textbf{Certified 3DEXPERIENCE Collaborative Industry Innovator}~\hfill {\color{emphasis}Dassault Systèmes}
  \item \textbf{MATLAB / Simulink Onramp}~\hfill {\color{emphasis}MathWorks}
\end{itemize}

\cvsection{ACTIVITÉS \& COMPÉTITIONS}
\begin{itemize}[leftmargin=*, label=\textbullet, nosep, itemsep=1.5pt]
  \item \textbf{Hack The Intelligence (AISEC 2026) - 1\textsuperscript{er} Prix :} Développement d’un prototype Smart Agriculture IoT.
  \item \textbf{Vision AI Competition (EMSI 2026) - 2\textsuperscript{ème} Prix :} Conception d’un système de traduction en langue des signes (ASL).
  \item \textbf{Congrès COCIA’26 - Présentation Scientifique :} Conception du prototype d’éclairage piézoélectrique SmartStruct.
  \item \textbf{Hackathon RallyIA (2026) - Participation :} Conception de l’architecture du projet "FindSpot" (Parking Intelligent).
\end{itemize}

\cvsection{LANGUES}
\textbf{Arabe :} Langue maternelle~\hfill \textbf{Français :} DELF B2~\hfill \textbf{Anglais :} EFSET B2

\end{document}
